
\section{Appendix: Existence of Regular Elements over Finite Fields}
\SGARefTarget{sga3:xiv:7}

\begin{center}
\itshape by J.-P. Serre
\end{center}

Throughout the Appendix, \(k\) denotes a finite field,
\(\bar k\) its algebraic closure, and
\(\mathcal G=\operatorname{Gal}(\bar k/k)\). If
\(q=\operatorname{Card}(k)\), then \(\mathcal G\) is
topologically generated by the Frobenius element
\[
  F\colon x\longmapsto x^q.
\]
We shall prove the following theorem.\footnote{The theorem is due
to Chevalley. The editor wishes to express his gratitude to
American Express which, by misplacing a trunk of Chevalley's
manuscripts, obliged him to reconstruct the proof.}

\paragraph{\textbf{Theorem.}}
\SGARefTarget{sga3:xiv:7:theorem}
{\itshape Let \(G\) be an adjoint semisimple group defined over
\(k\), and let \(\mathfrak g\) be its Lie algebra. Then the
\(k\)-Lie algebra \(\mathfrak g(k)\) contains a regular element.}

\paragraph{\textbf{Remarks.}}
\textup{\SGARefLink{sga3:xiv:7:lemma1}{(1)}} Recall that ``adjoint'' means that the center of
\(G\), as a subgroup scheme, is trivial. In view of the Chevalley
seminar, this also means that, for a maximal torus \(T\) of
\(G\), the character group \(X(T)\) of \(T_{\bar k}\) is
generated by the set \(R\) of roots. In particular, the rank of
\(\mathfrak g\) equals \(\dim T\), the rank of \(G\).

\textup{\SGARefLink{sga3:xiv:7:lemma2}{(2)}} The editor does not know whether the adjoint
hypothesis is indispensable.

\paragraph{\textbf{Lemma 1.}}
\SGARefTarget{sga3:xiv:7:lemma1}
{\itshape It is enough to prove the theorem when \(G\) is
geometrically simple.}

Here geometrically simple means that \(G_{\bar k}\) has no smooth
normal subgroup scheme other than \(G_{\bar k}\) and \((e)\);
equivalently, the associated root system is irreducible.

We may first suppose that \(G\) is indecomposable over \(k\).
Then \(G_{\bar k}\) is a product of geometrically simple
components \(S\), permuted transitively by \(\mathcal G\). If
\(\mathcal H\subset\mathcal G\) fixes one component \(S\), and
\(K\) is its fixed field, then \(S\) is defined over \(K\), and
the standard restriction-of-scalars argument gives
\[
  G=R_{K/k}(S).
\]
Similarly,
\[
  \mathfrak g=R_{K/k}(\mathfrak s),
\]
where \(\mathfrak s\) is the Lie algebra of \(S\). If the theorem
holds for \(S\), there is an element
\[
  x\in\mathfrak s(K)=\mathfrak g(k)
\]
regular in \(\mathfrak s\), and one checks immediately that it
is regular in \(\mathfrak g\). This proves the lemma.

Henceforth assume that \(G\) is geometrically simple, and choose
a maximal torus \(T\) of \(G\). Let \(X\) be the character group
of \(T_{\bar k}\), let \(R\) be its root system, \(W\) its Weyl
group, and \(E\) the group of automorphisms of \(X\) preserving
\(R\). Then \(W\) is normal in \(E\).

For another maximal torus \(T'\), write \(X',R',W',E'\) for the
corresponding objects. If
\(y\in G(\bar k)\) satisfies
\[
  yT_{\bar k}y^{-1}=T'_{\bar k},
\]
then \(\operatorname{int}(y)\) identifies
\((X',R',W',E')\) with \((X,R,W,E)\). Replacing \(y\) changes
this identification by an element of \(W\).

Frobenius acts on \(X'\) while preserving \(R'\), and thus
defines an element \(f_{T'}\in E'\). Put \(f=f_T\). After the
above identification \(E'\simeq E\), the element \(f_{T'}\)
becomes an element \(f'\in E\), well defined up to
\(W\)-conjugacy. We now compare \(f'\) with \(f\).
